\section{Statistiche d'Ordine e Problema della Mediana}

\subsection{Definizione}
\begin{definizione}[Statistica d'ordine]
Dato un insieme di $n$ elementi, la $i$-esima statistica d'ordine è l'$i$-esimo elemento più piccolo.
\end{definizione}

\begin{definizione}[Mediana]
Un valore $v$ è mediana se almeno il 50\% degli elementi è $\le v$ e almeno il 50\% è $\ge v$.
\begin{itemize}
	\item Se $n$ è dispari: mediana = statistica d'ordine $(n+1)/2$.
	\item Se $n$ è pari: mediana inferiore = statistica d'ordine $\lfloor (n+1)/2\rfloor$, mediana superiore = statistica d'ordine $\lceil (n+1)/2\rceil$.
\end{itemize}
Nelle slide, quando si scrive "mediana" si intende la \textbf{mediana inferiore}.
\end{definizione}

\begin{esempio}[Statistica d'ordine]
	Dato $A=\{3,5,9,4,6,1\}$ con $n=6$: la I statistica d'ordine è $1$ (minimo),
	la IV statistica d'ordine è $5$, la VI statistica d'ordine è $9$ (massimo).
\end{esempio}

\begin{esempio}[Mediana inferiore e superiore]
	Sempre con $A=\{3,5,9,4,6,1\}$: la mediana inferiore (III statistica d'ordine) è $4$;
	la mediana superiore (IV statistica d'ordine) è $5$.
\end{esempio}

\subsection{Intuizione}
Ordinare tutto l'array per trovare un solo rango è spesso inutile: il problema Selezione richiede un solo elemento, non la permutazione completa.

\subsection{Motivazione}
Selezione è un esempio didatticamente centrale di algoritmo divide et impera randomizzato con costo medio lineare.

\subsection{Algoritmo}
\begin{definizione}[Problema Selezione]
	\textbf{Input:} un insieme $A$ di $n$ numeri (distinti, per semplicità) e un intero $1\le i\le n$.\newline
	\textbf{Output:} l'$i$-esimo elemento più piccolo di $A$.
\end{definizione}

\paragraph{Minimo e massimo.}
\begin{algorithm}[H]
	\caption{\textsc{Minimum}$(A)$}
	\KwIn{Array $A[1:n]$}
	\KwOut{Valore minimo $m\in A$}
	$m\gets A[1]$\;
	\For{$i\gets 2$ \KwTo $n$}{
		\If{$m>A[i]$}{$m\gets A[i]$}
	}
	\Return $m$\;
\end{algorithm}
Richiede $n-1$ confronti: $\Theta(n)$, ottimale (ogni elemento non minimo deve confrontarsi almeno una volta).

\begin{algorithm}[H]
	\caption{\textsc{Maximum}$(A)$}
	\KwIn{Array $A[1:n]$}
	\KwOut{Valore massimo $M\in A$}
	$M\gets A[1]$\;
	\For{$i\gets 2$ \KwTo $n$}{
		\If{$M<A[i]$}{$M\gets A[i]$}
	}
	\Return $M$\;
\end{algorithm}
Analogo a \textsc{Minimum}, con $\Theta(n)$ confronti.

\begin{algorithm}[H]
	\caption{\textsc{MinMax}$(A)$}
	\KwIn{Array $A[1:n]$}
	\KwOut{Coppia $(m,M)$ minimo e massimo in $A$}
	\If{$n\bmod 2\neq 0$}{
		$m\gets M\gets A[1]$\;
		$i\gets 2$\;
	}
	\Else{
		\If{$A[1]<A[2]$}{$m\gets A[1]$;\; $M\gets A[2]$}
		\Else{$M\gets A[1]$;\; $m\gets A[2]$}
		$i\gets 3$\;
	}
	\While{$i\le n-1$}{
		\If{$A[i]<A[i+1]$}{
			\If{$m>A[i]$}{$m\gets A[i]$}
			\If{$M<A[i+1]$}{$M\gets A[i+1]$}
		}
		\Else{
			\If{$m>A[i+1]$}{$m\gets A[i+1]$}
			\If{$M<A[i]$}{$M\gets A[i]$}
		}
		$i\gets i+2$\;
	}
	\Return $(m,M)$\;
\end{algorithm}
Confronti: $\lceil\frac{3n}{2}\rceil-2$ (migliore di $2(n-1)$ di due scansioni separate).

\paragraph{Riduzione con pivot $v$.}
\[
A_L=\{x\in A\mid x<v\},\quad
A_v=\{x\in A\mid x=v\},\quad
A_R=\{x\in A\mid x>v\}.
\]
\[
\text{Selezione}(A,i)=
\begin{cases}
\text{Selezione}(A_L,i) & i\le |A_L|,\\
v & |A_L|<i\le |A_L|+|A_v|,\\
\text{Selezione}(A_R,i-|A_L|-|A_v|) & i>|A_L|+|A_v|.
\end{cases}
\]

\paragraph{Pseudocodice \textsc{RandSelect} (slide: CLRS 9.2).}
\begin{algorithm}[H]
	\caption{\textsc{RandSelect}$(A,p,r,i)$}
	\KwIn{Array $A$, indici $p,r$, rango locale $i$}
	\KwOut{$i$-esimo elemento più piccolo in $A[p:r]$}
	\If{$p=r$}{\Return $A[p]$\;}
	$q\gets\textsc{RandPartition}(A,p,r)$\;
	$k\gets q-p+1$\;
	\If{$i=k$}{\Return $A[q]$\;}
	\If{$i<k$}{
		\Return \textsc{RandSelect}$(A,p,q-1,i)$\;
	}
	\Else{
		\Return \textsc{RandSelect}$(A,q+1,r,i-k)$\;
	}
\end{algorithm}

\begin{algorithm}[H]
	\caption{\textsc{Partition}$(A,p,r)$}
	\KwIn{Array $A$, indici $p,r$; pivot in $A[r]$}
	\KwOut{Indice $q$ del pivot dopo la partizione}
	$x\gets A[r]$\;
	$i\gets p-1$\;
	\For{$j\gets p$ \KwTo $r-1$}{
		\If{$A[j]\le x$}{
			$i\gets i+1$\;
			scambia $A[i]$ con $A[j]$\;
		}
	}
	scambia $A[i+1]$ con $A[r]$\;
	\Return $i+1$\;
\end{algorithm}

\begin{algorithm}[H]
	\caption{\textsc{RandPartition}$(A,p,r)$}
	\KwIn{Array $A$, indici $p,r$}
	\KwOut{Indice $q$ restituito da \textsc{Partition}}
	$i\gets$ indice casuale in $\{p,\ldots,r\}$\;
	scambia $A[r]$ con $A[i]$\;
	\Return \textsc{Partition}$(A,p,r)$\;
\end{algorithm}

\subsection{Esempio guidato}
Sia $A=\{3,5,9,4,6,1\}$, $n=6$, mediana inferiore: $i=3$.
\begin{enumerate}
	\item Scelta casuale $v=3$:
	\[
	A_L=\{1\},\;A_v=\{3\},\;A_R=\{5,9,4,6\}.
	\]
	Poiché $3>|A_L|+|A_v|=2$, si passa a
	\[
	\text{Selezione}(A_R,3-1-1)=\text{Selezione}(\{5,9,4,6\},1).
	\]
	\item Scelta casuale $v=4$ nel nuovo insieme:
	\[
	A_L=\emptyset,\;A_v=\{4\},\;A_R=\{5,9,6\}.
	\]
	Poiché $0<1\le1$, la soluzione è $4$.
\end{enumerate}

\subsection{Grafico o diagramma}
\begin{center}
\begin{verbatim}
A = {3,5,9,4,6,1}, i=3
pivot=3
AL={1} | Av={3} | AR={5,9,4,6}
                 --> recurse on AR with i'=1

AR={5,9,4,6}, i'=1
pivot=4
AL={} | Av={4} | AR={5,9,6}
--> answer = 4
\end{verbatim}
\end{center}

\begin{figure}[H]
	\centering
	\begin{tikzpicture}[
		node distance=6mm and 8mm,
		box/.style={draw, primary, rounded corners, minimum width=1.6cm,
			minimum height=0.7cm, font=\small, align=center},
		arr/.style={-{Stealth[length=2mm]}, primary, thick}
		]
		\node[box] (A) {$A=\{3,5,9,4,6,1\}$\\$i=3$};
		\node[box, below=of A] (p) {pivot $v=3$};
		\node[box, below left=10mm and -5mm of p] (AL) {$A_L=\{1\}$};
		\node[box, below=of p] (Av) {$A_v=\{3\}$};
		\node[box, below right=10mm and -5mm of p] (AR) {$A_R=\{5,9,4,6\}$};
		\node[box, below=14mm of AR] (R1) {ricorsione\\$i'=1$};
		\draw[arr] (A) -- (p);
		\draw[arr] (p) -- (AL);
		\draw[arr] (p) -- (Av);
		\draw[arr] (p) -- (AR);
		\draw[arr] (AR) -- (R1) node[midway, right, font=\scriptsize] {$3>|A_L|+|A_v|$};
	\end{tikzpicture}
	\caption{Partizione ricorsiva di \textsc{RandSelect} sull'esempio della mediana inferiore.}
\end{figure}

\begin{figure}[H]
	\centering
	\begin{tikzpicture}[
		node distance=5mm and 6mm,
		box/.style={draw, primary, rounded corners, minimum width=1.4cm,
			minimum height=0.65cm, font=\small, align=center},
		arr/.style={-{Stealth[length=2mm]}, primary, thick}
		]
		\node[box] (AR) {$A_R=\{5,9,4,6\}$\\$i'=1$};
		\node[box, below=of AR] (p2) {pivot $v=4$};
		\node[box, below left=8mm and -2mm of p2] (AL2) {$A_L=\emptyset$};
		\node[box, below=of p2] (Av2) {$A_v=\{4\}$};
		\node[box, below right=8mm and -2mm of p2] (AR2) {$A_R=\{5,9,6\}$};
		\node[box, below=10mm of Av2, fill=cyan!10] (ans) {risposta $4$};
		\draw[arr] (AR) -- (p2);
		\draw[arr] (p2) -- (AL2);
		\draw[arr] (p2) -- (Av2);
		\draw[arr] (p2) -- (AR2);
		\draw[arr] (Av2) -- (ans) node[midway, right, font=\scriptsize] {$0<1\le 1$};
	\end{tikzpicture}
	\caption{Secondo passo ricorsivo dell'Esempio 3 delle slide.}
\end{figure}

\subsection{Dimostrazione}
\paragraph{Caso medio (come nelle slide).}
\begin{definizione}[Pivot buono]
	Un pivot $v$ è \textbf{buono} se cade tra il 25-esimo e il 75-esimo percentile di $A$.
	Il $j$-esimo \textbf{percentile} è il minimo valore sotto cui ricade almeno il $j\%$ degli elementi considerati.
\end{definizione}
Allora con probabilità $1/2$ il pivot è buono e il sottoproblema ha dimensione al più $3n/4$.

Indicando con $s(n)=\Theta(n)$ il costo di una \textsc{RandPartition} e con $w(n)$ il costo totale fino al primo pivot buono, le slide assumono:
\[
E[w(n)]=O(n).
\]
Quindi:
\[
E[T(n)]\le E[T(3n/4)]+O(n).
\]
Per sostituzione:
\[
E[T(n)]=O(n).
\]
Inoltre almeno una partizione completa è sempre necessaria, quindi:
\[
E[T(n)]=\Omega(n).
\]
Conclusione:
\[
E[T(n)]=\Theta(n).
\]

\begin{lemma}[Pivot buono e riduzione di dimensione]
	Se $A[q]$ è un pivot buono, allora $|A_L|\le\frac{3}{4}|A|$ e $|A_R|\le\frac{3}{4}|A|$.
\end{lemma}
\begin{proof}
	Per definizione di percentile, al più il 25\% degli elementi è strettamente minore di $v$ e al più il 25\% è strettamente maggiore; il resto appartiene ad $A_v$. Quindi $|A_L|,|A_R|\le n/4$ nel caso $|A_v|$ trascurabile; in generale $|A_L|+|A_v|+|A_R|=n$ implica $|A_L|,|A_R|\le\frac{3}{4}n$.
\end{proof}

\begin{softnote}
	Sia $X_v$ la variabile indicatrice dell'evento ``$v$ è un elemento buono''. Allora $\Pr\{X_v=1\}=E[X_v]=\frac{1}{2}$.
	Per due pivot casuali $A[q]$ e $A[q']$: $E[X_{A[q]}+X_{A[q']}]=1$, quindi il numero atteso di partizioni prima di un pivot buono è $2$ e $E[w(n)]=O(n)$.
\end{softnote}

\subsection{Analisi}
\paragraph{Caso peggiore.}
Se il pivot è sempre massimo (o minimo):
\[
T(n)=T(n-1)+\Theta(n)=\Theta(n^2).
\]

\paragraph{Caso medio.}
\[
E[T(n)]=\Theta(n).
\]

\paragraph{Nota implementativa dalle slide.}
\textsc{RandPartition} non usa memoria ausiliaria ma modifica l'ordine dell'array di input.

\subsection{Errori tipici d'esame}
\begin{hardnote}
Confondere "tempo medio" con "tempo nel caso peggiore": per \textsc{RandSelect} restano distinti.
\end{hardnote}

\begin{hardnote}
Dimenticare la correzione dell'indice nella chiamata su $A_R$:
$i' = i-|A_L|-|A_v|$.
\end{hardnote}

\subsection{Possibili domande d'esame}
\begin{itemize}
	\item Dimostrare che minimo/massimo richiedono almeno $n-1$ confronti.
	\item Completare \textsc{MinMax} con confronto a coppie e motivare il limite
	$\left\lceil\frac{3n}{2}\right\rceil-2$.
	\item Giustificare formalmente perché il caso medio di \textsc{RandSelect} è lineare.
	\item Riscrivere \textsc{RandSelect} con duplicati espliciti (partizione a tre insiemi).
\end{itemize}

% Median-of-Medians (\textsc{Select} deterministico, CLRS 9.3): fuori handout s05_1 (studio fino a 9.3 escluso).
